\section{Joint State, Event, and Source Problem}
\label{sec:problem-formulation}

\subsection{Hybrid DAE and Electrical Measurement Model}

Let $\bm x(t)\in\mathbb R^{n_x}$ denote differential states, $\bm z(t)\in\mathbb R^{n_z}$ algebraic states, $\bm u(t)$ scheduled inputs, $\bm q(t)$ network configuration, and $\bm\theta$ physical parameters. A physical intervention $\bm a_h(t)$ is indexed by hypothesis
\begin{equation}
 h=(c,\ell,\tau,\gamma,\bm d), \qquad h\in\mathcal H,
 \label{eq:physical-hypothesis}
\end{equation}
where $c$ is event family, $\ell$ is a bus or branch source, $\tau$ is onset, $\gamma$ is severity, and $\bm d$ describes the allowed temporal profile. The plant is
\begin{align}
 \dot{\bm x}(t)&=\bm f\!\left(\bm x,\bm z,\bm u,\bm q,\bm a_h;\bm\theta\right)+\bm w_x(t), \label{eq:dae-dynamic}\\
 \bm 0&=\bm g\!\left(\bm x,\bm z,\bm u,\bm q,\bm a_h;\bm\theta\right), \label{eq:dae-algebraic}\\
 \bm y_k&=\bm M_k\!\left[\bm h_{\mathcal P_k}(\bm x_k,\bm z_k,\bm q_k;\bm\theta)+\bm b_k\right]+\bm v_k. \label{eq:dae-measurement}
\end{align}
Here $\bm M_k$ selects the channels that are available at time $k$, $\bm b_k$ is a measurement-corruption process, and $\bm v_k$ is measurement noise. A dropout changes $\bm M_k$; bad data change $\bm b_k$; neither is a physical input to \eqref{eq:dae-dynamic}--\eqref{eq:dae-algebraic}. This separation permits a physical event and a measurement failure to coexist.

The algebraic network core is
\begin{equation}
 \bm I=\bm Y_{\mathrm{bus}}(\bm q,\bm\theta)\bm V,
 \qquad S_i=P_i+\mathrm jQ_i=V_i I_i^{*}. \label{eq:network-core}
\end{equation}
Writing $\bm Y_{\mathrm{bus}}=\bm G+\mathrm j\bm B$ and $V_i=|V_i|e^{\mathrm j\vartheta_i}$ gives
\begin{align}
 P_i&=|V_i|\sum_j|V_j|\!\left[G_{ij}\cos\vartheta_{ij}+B_{ij}\sin\vartheta_{ij}\right], \label{eq:ac-p}\\
 Q_i&=|V_i|\sum_j|V_j|\!\left[G_{ij}\sin\vartheta_{ij}-B_{ij}\cos\vartheta_{ij}\right], \label{eq:ac-q}
\end{align}
where $\vartheta_{ij}=\vartheta_i-\vartheta_j$. A line outage changes the branch primitives and therefore $\bm Y_{\mathrm{bus}}$; a bus fault adds a shunt-admittance intervention; generation and load changes modify the appropriate active/reactive injection equations. For a PMU at one terminal of branch $(i,j)$,
\begin{equation}
 I_{ij}=Y_{ff,ij}V_i+Y_{ft,ij}V_j,
 \qquad I_{ji}=Y_{tf,ij}V_i+Y_{tt,ij}V_j. \label{eq:terminal-current}
\end{equation}
The orientation, transformer tap, line charging, and shunt convention in \eqref{eq:terminal-current} are part of the inference contract; replacing terminal current with bus-injection current changes the likelihood.

\subsection{Sparse PMUs and Requested State Function}

Both evidence campaigns use PMU buses
\begin{equation}
 \mathcal P=\{2,5,6,10,19,22,29,39\},\qquad |\mathcal P|=8. \label{eq:pmu-set}
\end{equation}
The unobserved set is $\mathcal U=\{1,\ldots,39\}\setminus\mathcal P$, with $|\mathcal U|=31$. The state-estimation campaign stacks rectangular voltage and controlled terminal-current components,
\begin{equation}
 \bm y_{\mathcal P,k}=\operatorname{col}_{p\in\mathcal P}
 [\Re V_{p,k},\Im V_{p,k},\Re I_{p\ell,k},\Im I_{p\ell,k}], \label{eq:pmu-vector}
\end{equation}
and requests the hidden complex-voltage function
\begin{equation}
 \bm z_{\mathcal U,k}=\operatorname{col}_{i\in\mathcal U}[\Re V_{i,k},\Im V_{i,k}]\in\mathbb R^{62}. \label{eq:hidden-target}
\end{equation}
The event campaign also uses voltage angle, frequency, availability, and, in its longer-window implementation, current and ROCOF. Measurement semantics are kept campaign-specific rather than declaring channels equivalent across simulators.

\subsection{Eight Event Labels as Composable States}

Table~\ref{tab:event-ontology} defines the eight abnormal labels in the historical event bank. The scientific variables are not the label numbers. Classes 1--4 are single physical interventions; classes 5 and 7 are measurement-integrity states; classes 6 and 8 are compositions. Class 8 is a prescribed mixed event in the available simulations, not evidence of open-set recognition.

\begin{table*}[t]
\caption{Event ontology and historical ANDES inventory. Each source configuration has five replicas. ``Physical source'' and ``integrity source'' are separate outputs.}
\label{tab:event-ontology}
\centering
\footnotesize
\begin{tabularx}{\textwidth}{@{}cXlllr@{}}
\toprule
ID & Operational meaning & Physical intervention $c$ & Physical source $\ell$ & Integrity state/source & Records \\
\midrule
0 & Normal operation & none & none & nominal & 5 \\
1 & Bus fault & fault admittance & bus & nominal & 195 \\
2 & Line outage & topology removal & branch & nominal & 170 \\
3 & Generation change & active/reactive injection & generator bus & nominal & 50 \\
4 & Load change & active/reactive withdrawal & load bus & nominal & 90 \\
5 & Missing data & none & none & dropout / PMU & 40 \\
6 & Generation change with missing data & generation change & generator bus & dropout / PMU & 50 \\
7 & Bad data & none & none & corruption / PMU channel & 40 \\
8 & Prescribed mixed disturbance & generation plus load changes & declared generator bus & measurement disturbance / PMU & 50 \\
\midrule
 & Total & & & & 690 \\
\bottomrule
\end{tabularx}
\end{table*}

The internal decision at time $k$ is
\begin{equation}
 \left(\widehat{\bm z}_{\mathcal U,k},\bm\Sigma_{z,k},
 \widehat c_k,\widehat\ell^{\mathrm{phys}}_k,
 \widehat m_k,\widehat\ell^{\mathrm{meas}}_k,\mathcal C_k\right),
 \label{eq:joint-output}
\end{equation}
where $m_k$ denotes the integrity state and $\mathcal C_k\subseteq\mathcal H$ is a credible candidate set. A location is empty when it does not apply. The candidate set is required because limited observability can make a forced Top-1 decision scientifically unjustified.

\subsection{Correctness and Evaluation Units}

State accuracy uses hidden-bus total vector error
\begin{equation}
 \operatorname{TVE}^{(r)}=\frac{1}{T_r|\mathcal U|}
 \sum_{k=1}^{T_r}\sum_{i\in\mathcal U}
 \frac{|\widehat V_{i,k}-V_{i,k}|}{\max(|V_{i,k}|,\epsilon)}, \label{eq:tve}
\end{equation}
reported as a percentage. Event detection, event-family macro-F1, exact physical Top-1/Top-3, and integrity-source Top-1 use their own declared populations. The joint correctness indicator is one only when every applicable component of \eqref{eq:joint-output} is correct; a correct event family cannot conceal a wrong source. Confidence intervals use an independently generated trajectory, event parent, or rebuilt plant as the statistical unit, never individual 30-Hz samples.

The causal contract permits only $\bm y_{0:k}$ at time $k$. A fixed-lag smoother, complete-window statistic, or centered temporal feature is retrospective and is labeled accordingly. Source generalization requires that physical source $\ell$ be absent from all source-labeled training and tuning, not merely that the test trajectory use a new random seed.
